\section{Detailed Implementation Protocol: The ``Staged Gate'' Workflow}

This section synthesizes the research into a coherent, actionable workflow for the research team.

\subsection{Phase 1: Data Ingestion and Static Auditing (CPU)}
This phase runs on the entire generated dataset.

\begin{enumerate}
    \item \textbf{Load and Normalize:}
    \begin{itemize}
        \item Load $(\vect{q}, \vect{p}, t)$ from Brutus output.
        \item Check for NaN or Inf.
        \item Compute $SCS$ (Energy Drift). \textbf{Gate:} Reject if $|SCS| > 10^{-5}$.
    \end{itemize}
    \item \textbf{Stiffness Analysis:}
    \begin{itemize}
        \item Compute accelerations $\vect{a}_i$ and velocities $\vect{v}_i$.
        \item Compute $SASD$ profile.
        \item \textbf{Gate:} If $\max(SASD) > 0.5$, flag for ``Resampling.'' Do not train on this trajectory yet. It requires finer Brutus output steps.
    \end{itemize}
    \item \textbf{Manifold Check (IPR):}
    \begin{itemize}
        \item Fit Cubic Hermite Spline.
        \item Sample 10 synthetic points per interval.
        \item Compute $IPR$.
        \item \textbf{Gate:} If $IPR > 10^{-3}$, the data is aliased. Reject or Resample.
    \end{itemize}
\end{enumerate}

\subsection{Phase 2: The ``Overfit Challenge'' (GPU)}
This phase runs on a randomly selected \textit{subset} (e.g., 1\%) of the passing data to validate the configuration.

\begin{enumerate}
    \item \textbf{Setup:} Initialize the exact CTL-PINN architecture intended for full training.
    \item \textbf{Challenge:} Attempt to overfit a \textit{single} stiff trajectory (one that passed Phase 1 but had high SASD).
    \item \textbf{Monitoring:}
    \begin{itemize}
        \item Track Loss $\mathcal{L}_{train}$.
        \item Track Causal Weights $w(t)$ evolution.
        \item Track Gradient Norms $\norm{\nabla_\theta \mathcal{L}}$.
    \end{itemize}
    \item \textbf{Success Criteria:}
    \begin{itemize}
        \item $\mathcal{L}_{train}$ converges to machine precision ($< 10^{-6}$).
        \item Causal weights propagate to $t=T$ (i.e., the ``firewall'' is breached).
    \end{itemize}
    \item \textbf{Diagnostic:}
    \begin{itemize}
        \item If loss plateaus: Architecture is too small (under-parameterized).
        \item If weights don't propagate: Causal parameter $\lambda$ is too high (too strict).
        \item If gradients explode: Normalization is incorrect.
    \end{itemize}
\end{enumerate}

\subsection{Phase 3: Dataset Composition and Augmentation}

\begin{enumerate}
    \item \textbf{Diversity Check:} Compute Phase Space Entropy ($H_{data}$) on the accepted dataset.
    \item \textbf{Balancing:} If the dataset is dominated by stable orbits (low entropy), use \textbf{Importance Resampling}. Run Brutus specifically initialized in ``chaotic'' regions (e.g., using the figure-8 or Pythagorean 3-body initial conditions) to generate more high-entropy examples~\cite{ref31}.
    \item \textbf{Symmetry Augmentation:} Apply $SO(3)$ rotations and translations to valid trajectories to increase dataset size without computational cost, ensuring the PINN learns rotational invariance~\cite{ref6}.
\end{enumerate}